% ============================================================
% sections/sixteen_atom_closure.tex
% "The sixteen-atom closure (t = 4)"
% for problems/23/writeup/arxiv/rotor_window_closures
%
% Conventions match the companion paper
% problems/23/writeup/arxiv/shortest_support_obstructions/main.tex:
% 11pt article; amsmath, amssymb, amsthm, booktabs, enumitem;
% plain-style theorem/proposition/lemma/corollary and
% definition-style definition/remark, numbered by section;
% macros \bip, \supp, \mc as in the companion.
%
% Bibliography keys used by this section (assembler: provide these):
%   SSO            the companion paper (shortest-geodesic supports in
%                  triangle-free maximum cuts, arXiv 2026)
%   McKayPiperno14 B. D. McKay and A. Piperno, Practical graph
%                  isomorphism II, J. Symbolic Comput. 60 (2014), 94-112
%   Mantel07       W. Mantel, Problem 28, Wiskundige Opgaven 10 (1907),
%                  60-61
% ============================================================

\section{The sixteen-atom closure}\label{sec:t4}

The classification of the companion paper \cite{Ferudun26supports} stops at ten atoms.
This section treats the first structured window beyond it: support
circuits with sixteen atoms carrying two \emph{rotating owners} of width
four. We first prove, by hand, the crossover bounds that isolate this
window and leave it a margin of exactly one support edge; we then close
the window by an exhaustive computation with two independent verifiers.
The closure eliminates one local repair mechanism for Hall expansion.
It does not resolve the Erdos \(n^2/25\) conjecture.

\subsection{Support circuits}\label{subsec:t4:circuits}

Throughout this section \(F\) is a connected bipartite graph with parts
\(X\) and \(Y\), and \(\mathcal A\) is a set of \(m\) distinct
\emph{atoms}: pairs \(\{p,q\}\) of vertices of \(F\) with
\(d_F(p,q)=4\). Both endpoints of an atom lie in one part. A \emph{row}
of an atom \(a\) is a four-edge \(F\)-geodesic between its endpoints,
written as a vertex sequence \((p_0,p_1,p_2,p_3,p_4)\); the support
\(P_F(a)\) is the union of the edge sets of all rows of \(a\). An
\(m\)-footprint obstruction is defined as in the companion paper: \(F\)
has \(m-1\) edges, the graph \((V(F),\mathcal A)\) is triangle-free, the
supports cover \(E(F)\), and every edge of \(F\) lies in at least two
supports.

\begin{definition}\label{def:t4:circuit}
A \emph{support circuit} is an \(m\)-footprint obstruction
\((F,\mathcal A)\) such that, in addition, for every \(a\in\mathcal A\)
the family \(\{P_F(b):b\in\mathcal A\setminus\{a\}\}\) admits a system
of distinct representatives that uses every edge of \(F\) exactly once.
\end{definition}

The name records that \(\mathcal A\) is a circuit of the transversal
matroid of the support incidence system: deleting any one atom restores
a perfect matching onto the support edges.

\begin{lemma}[Transfer]\label{lem:t4:transfer}
Let \(G\) be triangle-free with a cut \((V_0,V_1)\), crossing graph
\(B\) connected, and let \(\mathcal S\subseteq M\) be inclusion-minimal
with \(|\supp_B(\mathcal S)|<|\mathcal S|=m\). Put
\(F=\supp_B(\mathcal S)\). Then \((F,\mathcal S)\) is an \(m\)-atom
support circuit, and its atom--edge incidence graph is connected.
\end{lemma}

\begin{proof}
The minimal-footprint lemma of the companion paper \cite{Ferudun26supports} gives all
conditions of an \(m\)-footprint obstruction, with every atom at
\(F\)-distance four and its support equal to the union of its
length-four \(F\)-geodesics. For the representatives, fix
\(a\in\mathcal S\). By minimality every subfamily of
\(\mathcal S\setminus\{a\}\) satisfies Hall's condition inside \(F\), so
by Hall's theorem \(\mathcal S\setminus\{a\}\) has a system of distinct
representatives in \(E(F)\); since
\(|\mathcal S\setminus\{a\}|=m-1=|E(F)|\), the system uses every edge.
For connectivity, split the incidence graph into components with atom
classes \(\mathcal A_i\) and edge classes \(F_i\). Every edge lies in
some support, so every component contains an atom, and every support is
contained in the component of its atom. If there were at least two
components, then from \(\sum_i|F_i|=m-1<\sum_i|\mathcal A_i|\) some
component would satisfy \(|F_i|<|\mathcal A_i|\), and \(\mathcal A_i\)
would be a proper support-deficient subfamily, contradicting
minimality.
\end{proof}

\begin{lemma}\label{lem:t4:induced}
Every row is an induced path of \(F\): if \(u,w\) lie on a row \(R\)
and \(uw\in E(F)\), then \(uw\in E(R)\).
\end{lemma}

\begin{proof}
A chord between positions \(i<j\) of \(R\), \(j-i\ge2\), yields a walk
of length \(i+1+(4-j)\le3\) between the endpoints of \(R\)'s atom,
contradicting distance four.
\end{proof}

\begin{definition}\label{def:t4:selection}
A \emph{selection} \(\sigma\) chooses one row \(\sigma(a)\) for each
atom \(a\in\mathcal A\); its \emph{selected support}
\(\supp(\sigma)\) is the union of the edge sets of the chosen rows.
Given \(\sigma\), the star of a vertex \(v\) is
\emph{\(\sigma\)-covered} if there is a neighbour \(x_0\) of \(v\)
in \(F\) (the \emph{active} neighbour) such that
\(vx_0\notin\supp(\sigma)\) and, for every other \(F\)-neighbour
\(y\) of \(v\), some chosen row contains both \(x_0\) and \(y\).
\end{definition}

\begin{definition}\label{def:t4:rotating}
Let \((F,\mathcal A)\) be a support circuit. A \emph{rotating family of
width \(t\)} is a set of \(k\) distinct vertices
\(v_1,\dots,v_k\), all in one part of \(F\), such that every \(v_i\)
has at least \(t\) neighbours in \(F\) and is an endpoint of at least
\(t\) atoms, and every two \(v_i,v_j\) have a common \(F\)-neighbour.
Its members are called \emph{owners}.
\end{definition}

The window \((t,k)\) arises from a local strategy that tries to repair
failed Hall expansion by re-routing supports around \(k\) designated
vertices whose stars rotate between selections. Whatever the merits of
that strategy, its combinatorial footprint is precise: a
\(t^2\)-atom support circuit carrying a rotating family of \(k\) owners
of width \(t\) whose stars are covered by a common selection. The
statements below concern only this footprint.

\subsection{Crossover bounds}\label{subsec:t4:crossover}

\begin{theorem}\label{thm:t4:crossover}
Let \((F,\mathcal A)\) be an \(m\)-atom support circuit carrying a
rotating family \(v_1,\dots,v_k\) of width \(t\ge2\).
\begin{enumerate}[label=\textup{(\roman*)}]
\item \(|E(F)|\ge kt+t\); equivalently \(m\ge kt+t+1\).
\item If \(k=2\) and the star of \(v_1\) is \(\sigma\)-covered for some
selection \(\sigma\), then \(|E(F)|\ge3t+2\); equivalently
\(m\ge3t+3\).
\end{enumerate}
\end{theorem}

\begin{proof}
Say the owners lie in \(X\).

(i) Distinct vertices of one part are never adjacent in \(F\), so the
\(F\)-stars of the owners are pairwise disjoint edge sets, and they
contribute at least \(kt\) edges. Now fix \(v=v_1\) and distinct
atom-partners \(b_1,\dots,b_t\) of \(v\); each satisfies
\(d_F(v,b_j)=4\), so \(b_j\) is not an owner: an owner \(v_l\ne v\) has
a common neighbour with \(v\) and hence \(d_F(v,v_l)=2\). For each
\(j\) choose a row \(R_j\) of the atom \(\{v,b_j\}\) and let \(f_j\) be
its terminal edge at \(b_j\), with endpoints \(b_j\in X\) and
\(q_j\in Y\). The \(f_j\) are pairwise distinct, since their
\(X\)-endpoints \(b_j\) are distinct; and none is incident with an
owner, since \(q_j\in Y\) while the owners lie in \(X\), and
\(b_j\) is not an owner. This gives \(t\) further edges, hence
\(|E(F)|\ge kt+t\), and \(m-1=|E(F)|\) finishes (i).

(ii) Keep the \(2t\) owner-incident edges and the \(t\) terminal edges
from (i), taken at \(v=v_1\). Let \(x_0\) be the active neighbour and
let \(y\ne x_0\) be any other \(F\)-neighbour of \(v\) (it exists since
\(t\ge2\)). By coverage some chosen row \(Q\) contains \(x_0\) and
\(y\). First, \(v\notin Q\): otherwise \(v\) and \(x_0\) both lie on
the induced path \(Q\) (Lemma~\ref{lem:t4:induced}) and
\(vx_0\in E(F)\), so \(vx_0\) would be an edge of \(Q\subseteq
\supp(\sigma)\), contradicting the choice of \(x_0\). Hence no edge of
\(Q\) is incident with \(v\), and at most two of the four edges of
\(Q\) are incident with \(v_2\), so at least two edges of \(Q\) are
incident with no owner.

It remains to show that no edge of \(Q\) equals a terminal edge
\(f_j\). Suppose \(f_j\in E(Q)\); then \(b_j\in V(Q)\). The vertices
\(x_0,y\) are \(F\)-neighbours of \(v\in X\), so \(x_0,y\in Y\), while
\(b_j\in X\). If the atom of \(Q\) lies in \(X\), the \(Y\)-vertices of
\(Q\) are its two internal odd positions, so \(\{x_0,y\}\) is exactly
that pair, and every \(X\)-vertex of \(Q\) --- in particular \(b_j\)
--- has all its \(Q\)-neighbours in \(\{x_0,y\}\). If the atom of
\(Q\) lies in \(Y\), then \(b_j\) occupies an internal odd position and
its two \(Q\)-neighbours lie in the three \(Y\)-positions, which
include \(x_0\) and \(y\); by pigeonhole \(b_j\) is again
\(Q\)-adjacent to \(x_0\) or \(y\). Either way \(F\) contains an edge
from \(b_j\) to a neighbour of \(v\), so \(d_F(v,b_j)\le2\),
contradicting \(d_F(v,b_j)=4\). Thus the two owner-free edges of \(Q\)
are new, and \(|E(F)|\ge2t+t+2=3t+2\).
\end{proof}

\begin{corollary}[Crossover table]\label{cor:t4:table}
Let \((F,\mathcal A)\) be a support circuit with \(m=t^2\) atoms,
\(t\ge3\), carrying a rotating family of \(k\ge2\) owners of width
\(t\) whose stars include one covered star when \(k=2\). Then:
\begin{enumerate}[label=\textup{(\alph*)}]
\item \(k\ge t-1\) is impossible; in particular no such family exists
for \(t=3\);
\item for \(k=2\), necessarily \(t\ge4\);
\item at \(t=4\) the only surviving pair is \((t,k)=(4,2)\), and its
margin is a single edge: the bound \(3t+2=14\) against the budget
\(t^2-1=15\).
\end{enumerate}
\end{corollary}

\begin{proof}
(a) If \(k\ge t-1\) then \(kt+t\ge t^2>t^2-1=|E(F)|\), contradicting
Theorem~\ref{thm:t4:crossover}(i); for \(t=3\) this excludes every
\(k\ge2\). (b) \(3t+2\le t^2-1\) fails for \(t=3\) (\(11>8\)) and holds
first at \(t=4\) (\(14\le15\)). (c) At \(t=4\), part (a) kills
\(k\ge3\) (\(kt+t\ge16>15\)).
\end{proof}

\begin{table}[ht]
\centering
\begin{tabular}{c@{\quad}c@{\quad}c@{\quad}c}
\toprule
\(t\) & budget \(t^2-1\) & \(k=2\) bound \(3t+2\) & \(k=3\) bound \(4t\)\\
\midrule
3 & 8  & 11 (excluded) & 12 (excluded)\\
4 & 15 & 14 (margin 1) & 16 (excluded)\\
5 & 24 & 17 (margin 7) & 20 (margin 4)\\
6 & 35 & 20 (margin 15) & 24 (margin 11)\\
\bottomrule
\end{tabular}
\caption{The crossover bounds of Theorem~\ref{thm:t4:crossover} against
the edge budget of a \(t^2\)-atom support circuit. All families with
\(k\ge t-1\) are excluded outright. The first surviving window is
\((t,k)=(4,2)\), with margin exactly one.}
\label{tab:t4:crossover}
\end{table}

\subsection{The two-owner window at \texorpdfstring{\(t=4\)}{t=4}}
\label{subsec:t4:window}

\begin{definition}\label{def:t4:window}
A \emph{two-owner window} is a triple \((F,\mathcal A,\{v,m\})\), where
\((F,\mathcal A)\) is a \(16\)-atom support circuit and \(v\ne m\) are
vertices in one part of \(F\) such that
\(\deg_F(v)=\deg_F(m)=4\), \(v\) and \(m\) have at least two common
\(F\)-neighbours, and each of \(v,m\) is an endpoint of exactly four
atoms. The window is \emph{covered} if a single selection
\(\sigma\)-covers the stars of both owners.
\end{definition}

The exact degree and atom counts in Definition~\ref{def:t4:window}
reflect the graph-level states behind the window, in which each owner
has exactly four crossing and four monochromatic neighbours. That the
\emph{full} crossing star of a covered owner is forced into the support
union --- so that \(\deg_F=4\) rather than \(\deg_F\le4\) --- is the
content of the next lemma.

\begin{lemma}[Star absorption]\label{lem:t4:absorption}
In the setting of Lemma~\textup{\ref{lem:t4:transfer}}, let \(v\) be a
vertex with exactly four \(B\)-neighbours, let \(\sigma\) select one
\(B\)-geodesic row for each atom of \(\mathcal S\), and suppose there
is a \(B\)-neighbour \(x_0\) of \(v\) with
\(vx_0\notin\supp(\sigma)\) such that each pair \(\{x_0,y\}\), for the
other three \(B\)-neighbours \(y\) of \(v\), lies on a common chosen
row. Then all four \(B\)-star edges of \(v\) belong to
\(F=\supp_B(\mathcal S)\), and \(\deg_F(v)=4\).
\end{lemma}

\begin{proof}
Fix \(y\) and a chosen row \(Q\) containing \(x_0\) and \(y\). As in
Theorem~\ref{thm:t4:crossover}(ii), \(v\notin Q\): rows are induced
(Lemma~\ref{lem:t4:induced} applies verbatim to \(B\)-geodesics), and
\(v\in Q\) would force the non-selected edge \(vx_0\) into \(Q\). The
vertices \(x_0,y\) lie in the shore opposite \(v\), hence at even
distance along \(Q\); distance four would make \(\{x_0,y\}\) the atom
of \(Q\), impossible because \(d_B(x_0,y)=2\) through \(v\). So
\(x_0,y\) lie at \(Q\)-distance two, with a middle vertex
\(q\ne v\). Replacing \(q\) by \(v\) turns \(Q\) into another
four-edge \(B\)-path with the same endpoints --- a path because
\(v\notin Q\) --- hence another row of the same atom. Its edges
\(vx_0\) and \(vy\) therefore lie in the complete support of that atom,
which is contained in \(F\).
\end{proof}

\begin{proposition}[The \(K_{2,3}\) core]\label{prop:t4:core}
In every two-owner window \((F,\mathcal A,\{v,m\})\) the owners share
an atom-partner: there is a vertex \(b\) with
\(\{v,b\},\{m,b\}\in\mathcal A\). Consequently, if \(x,y\) are two
common neighbours of the owners, the window contains the mixed
\(K_{2,3}\) core on \(\{v,m\}\cup\{x,y,b\}\): the four support edges
\(vx,vy,mx,my\) and the two atoms \(\{v,b\},\{m,b\}\).
\end{proposition}

\begin{proof}
The two owner stars are disjoint and contribute eight of the fifteen
edges. As in Theorem~\ref{thm:t4:crossover}(i), each owner has four
distinct terminal edges, none incident with an owner; so the eight star
edges and the terminal edges of \(v\) and of \(m\) would number
\(8+4+4=16>15\) if the two terminal families were disjoint. Hence some
terminal edge of \(v\) equals one of \(m\); equal edges have equal
endpoints in the owners' part, so the corresponding atom-partners
coincide.
\end{proof}

The finite incidence core of Proposition~\ref{prop:t4:core} is
kernel-checked in Lean~4 as
\texttt{K2TailIncidence.exists\_common\_tail\_of\_support\_card\_le\_fifteen}
(module \texttt{R44K2TailOverlap.lean}, source SHA-256 prefix
\texttt{12DFB927}), compiling with axioms exactly \texttt{propext},
\texttt{Classical.choice}, \texttt{Quot.sound}. It is a structural
reduction, not by itself an exclusion.

\begin{lemma}[Vertex range]\label{lem:t4:vertexrange}
Every two-owner window satisfies \(8\le|V(F)|\le15\). If some
selection covers the star of one owner, then \(|V(F)|\le14\).
\end{lemma}

\begin{proof}
Lower bound: \((V(F),\mathcal A)\) is triangle-free with \(16\) edges,
so \(16\le|V(F)|^2/4\) by Mantel's theorem \cite{Mantel07}, giving
\(|V(F)|\ge8\). Upper bound: \(F\) is connected with \(15\) edges, so
\(|V(F)|\le16\) with equality only for a tree; but two common
neighbours \(x,y\) of the owners give the four-cycle \(v\)-\(x\)-\(m\)-\(y\),
so \(F\) is not a tree and \(|V(F)|\le15\).

Now suppose \(|V(F)|=15\) and the star of \(v\) is \(\sigma\)-covered,
with active neighbour \(x_0\) and other neighbours \(y_1,y_2,y_3\).
A connected graph with equally many vertices and edges contains exactly
one cycle. For each \(i\) take a chosen row \(Q_i\) containing \(x_0\)
and \(y_i\). Exactly as in Lemma~\ref{lem:t4:absorption} (whose
argument is purely about \(F\)-geodesics): \(v\notin Q_i\); the pair
\(x_0,y_i\) cannot lie at \(Q_i\)-distance four, since that would make
\(\{x_0,y_i\}\) the atom of \(Q_i\) while \(v\) is a common
\(F\)-neighbour, so \(d_F(x_0,y_i)=2\ne4\); hence \(x_0,y_i\) lie at
\(Q_i\)-distance two with middle \(q_i\ne v\). Then
\[
 C_i\;=\;v\hbox{-}x_0\hbox{-}q_i\hbox{-}y_i\hbox{-}v
\]
is a four-cycle of \(F\), using the star edges \(vx_0,vy_i\) and two
edges of \(Q_i\). For \(i\ne j\) the cycles differ: \(y_i\in C_i\),
while \(y_i\notin\{v,x_0,q_j,y_j\}\) (the middles \(q_j\) lie in the
owners' part, the \(y\)'s in the other). Three distinct cycles
contradict unicyclicity, so \(|V(F)|\le14\).
\end{proof}

\subsection{The census}\label{subsec:t4:census}

The closure now rests on one exhaustive computation.

\begin{theorem}[Computer-assisted]\label{thm:t4:census}
Up to isomorphism there are exactly \(153{,}978\) connected bipartite
graphs with \(15\) edges and between \(8\) and \(15\) vertices
\textup{(}Table~\textup{\ref{tab:t4:geng}}\textup{)}. Of these, exactly
\(34\) admit two same-part degree-four vertices \(v,m\) with at least
two common neighbours, at least four same-part vertices at distance
four from each, and a common same-part vertex at distance four from
both; each of the \(34\) admits exactly one such pair. These
embeddings extend to exactly \(576\) two-owner windows
\((F,\mathcal A,\{v,m\})\), through \(74{,}920\) atom completions, of
which \(2{,}299\) have a triangle-free atom graph, \(862\) satisfy
support coverage and edge multiplicity at least two, and \(576\)
satisfy all sixteen deletion representative systems. Furthermore:
\begin{enumerate}[label=\textup{(\roman*)}]
\item all \(576\) windows have \(|V(F)|=15\), on exactly four
isomorphism types of \(F\); in particular no two-owner window has at
most \(14\) support vertices;
\item in every window, at least eight atoms have every row through
\(v\), and at least eight have every row through \(m\); over the
\(576\) windows the number of such atoms has distribution
\(\{8{:}255,\ 9{:}193,\ 10{:}101,\ 11{:}26,\ 12{:}1\}\), identically
for \(v\) and \(m\);
\item no complete row family of any atom of any window contains a
\emph{middle swap}: two rows \((a,x,v,y,b)\) and \((a,x,m,y,b)\)
differing precisely in the middle vertex, one through each owner.
\end{enumerate}
\end{theorem}

The reduction to a finite computation is complete for two-owner
windows: \(F\) is connected, bipartite, with fifteen edges and, by
Lemma~\ref{lem:t4:vertexrange}, between \(8\) and \(15\) vertices; the
degree, common-neighbour and atom-count conditions are part of
Definition~\ref{def:t4:window}; the four atom-partners of each owner
are same-part vertices at distance four; and the common distance-four
vertex exists by Proposition~\ref{prop:t4:core}. The computation
generated all candidate graphs with \texttt{geng} from nauty/Traces
\cite{McKayPiperno14} (\texttt{geng -c -b \(n\) 15:15},
\(8\le n\le15\)), located the owner embeddings (the filter is invariant
under swapping the two colour classes), enumerated all admissible
sixteen-atom completions, and tested the circuit conditions with exact
integer arithmetic. Every
complete row family of every surviving window was then reconstructed
and every pair of rows in every family was tested for the middle-swap
pattern; the row-selection products of the \(576\) windows comprise
\(16{,}288\) selections in total.

\begin{table}[ht]
\centering
\begin{tabular}{c@{\quad}r@{\qquad}c@{\quad}r}
\toprule
\(n\) & graphs & \(n\) & graphs\\
\midrule
8  & 2      & 12 & 15{,}285\\
9  & 30     & 13 & 36{,}337\\
10 & 496    & 14 & 52{,}909\\
11 & 3{,}675 & 15 & 45{,}244\\
\bottomrule
\end{tabular}
\caption{Connected bipartite graphs with \(15\) edges, by vertex count:
the census ground set, \(153{,}978\) graphs in total. The counts were
produced by \texttt{geng}, reproduced independently by the
NetworkX-based verifier, and replayed once more for this paper.}
\label{tab:t4:geng}
\end{table}

\begin{corollary}[The sixteen-atom closure]\label{cor:t4:closure}
There is no covered two-owner window. Consequently, in a triangle-free
graph with a cut whose crossing graph is connected, no inclusion-minimal
support-deficient family of sixteen bad edges admits two same-shore
vertices of crossing degree four, with at least two common crossing
neighbours and exactly four atoms each, whose stars are simultaneously covered
by one selection of shortest rows. The window \((t,k)=(4,2)\) of
Table~\textup{\ref{tab:t4:crossover}} is empty.
\end{corollary}

\begin{proof}
Suppose \((F,\mathcal A,\{v,m\})\) is a covered two-owner window. By
Lemma~\ref{lem:t4:vertexrange} it has at most \(14\) support vertices;
by Theorem~\ref{thm:t4:census}(i) no two-owner window --- covered or
not --- has fewer than \(15\). For the graph-level statement, transfer
by Lemma~\ref{lem:t4:transfer}: the minimal family induces a
\(16\)-atom support circuit whose atoms, rows and supports coincide
with the graph-level ones, and Lemma~\ref{lem:t4:absorption} forces
each covered owner's full crossing star into \(F\), so
\(\deg_F(v)=\deg_F(m)=4\) and the selection covers both stars inside
\((F,\mathcal A)\).
\end{proof}

Two features of Theorem~\ref{thm:t4:census} make the closure robust.
First, parts (ii) and (iii) exclude the window a second time, at
\(|V(F)|=15\), independently of the covered-star analysis: the
re-routing mechanism that motivates the window both prescribes a
routing load of four atoms through each owner --- against the forced
load of at least eight in (ii) --- and requires at least one middle
swap in some complete row family, against (iii). Part (iii) is
quantified over every pair of rows of every family of every window,
with no side conditions. Second, the middle-swap exclusion has a
row-free confirmation:

\begin{proposition}\label{prop:t4:swapgeometry}
If two rows of one atom form a middle swap \((a,x,v,y,b)\),
\((a,x,m,y,b)\), then \(x,y\) are distinct common neighbours of \(v\)
and \(m\), \(a\) is a neighbour of \(x\) outside \(\{v,m\}\), \(b\) is
a neighbour of \(y\) outside \(\{v,m\}\), and \(\{a,b\}\) is a
distance-four pair witnessed by both displayed paths. This implication
is kernel-checked in Lean~4
\textup{(}\texttt{LiveMiddleSwapCrossOuter.lean}, theorem
\texttt{live\_middle\_swap\_has\_cross\_outer}, source SHA-256 prefix
\texttt{3DFF7897}; axioms exactly \texttt{propext},
\texttt{Quot.sound}, attested by both the acceptance record and a
fresh rebuild-and-probe at assembly,
\texttt{anc/lean\_axiom\_probe/}\textup{)}. Exhausting all ordered choices of
\((x,y,a,b)\) over the four support types of the \(576\) windows
\textup{(}multiplicities \(180,190,190,16\)\textup{)} yields zero such
configurations \textup{(}artifact SHA-256 prefix \texttt{79db75b9},
verdict \texttt{PASS\_NO\_LIVE\_MIDDLE\_SWAP\_GEOMETRY}\textup{)}, so
Theorem~\textup{\ref{thm:t4:census}(iii)} follows again without
enumerating rows.
\end{proposition}

\begin{remark}[An independent catalogue and a tight near-miss]
\label{rem:t4:nearmiss}
A second, independently designed session enumeration attacked the same
window by vertex count and reported: no owner distance-four profile at
\(|V(F)|\in\{10,11\}\); at most ten same-part distance-four pairs at
\(|V(F)|=12\) (short of sixteen); \(280\) candidate graphs at
\(|V(F)|=13\) with no sixteen-atom circuit carrying four owner-star
atoms at each owner; and \(455\) candidate graphs at \(|V(F)|=14\) with
no simultaneously covered owner stars. In that route, distance-four
pairs within one part are certified by the exact criterion
\((Q_L)_{uv}=0<(Q_L^2)_{uv}\), where \(Q_L=A_HA_H^{\mathsf T}\) and
\(A_H\) is the biadjacency matrix. This catalogue is a single
implementation and we use it only as corroboration; the proof of
Corollary~\ref{cor:t4:closure} rests on the dual-verified census. The
catalogue also exhibited a tight near-miss at \(|V(F)|=14\) (checker
SHA-256 prefix \texttt{4644e5ab}), whose archived invariants we record:
parts \(\{0,\dots,6\}\) and \(\{7,\dots,13\}\); fifteen support edges
and sixteen atoms; owners \(0\) and \(1\) sharing the two support
neighbours \(7,8\) and the three atom-partners \(3,5,6\); shared
terminal edges \(\{3,13\},\{5,13\},\{6,13\}\); the remaining edge
budget spent on \(\{2,11\}\) and \(\{4,9\}\); a complete row database
of \(864\) selections, none of which covers both owner stars. The
recorded failure is local: the common neighbours \(7,8\) lie on no
geodesic row of the relevant families, and every alternative row
passes through an owner and thus absorbs the would-be active edge.
The full edge and atom lists of this configuration were
session-internal and are not reproduced here.
\end{remark}

\begin{remark}[Why geometry is load-bearing]\label{rem:t4:abstract}
The census cannot be replaced by support-cardinality reasoning. There
is an abstract sixteen-atom, fifteen-edge transversal circuit with
four-element supports --- artifact SHA-256 prefix \texttt{5b386cd9},
independently re-verified --- realizing the forced tail overlap of
Proposition~\ref{prop:t4:core}, edge multiplicity at least two, full
deficiency one, and Hall's condition on all \(65{,}534\) proper atom
subsets. Its supports have not been realized as geodesics of any one
graph, and no census window realizes them.
Any proof of Corollary~\ref{cor:t4:closure} must therefore use the
geodesic geometry, as the census does.
\end{remark}

\subsection*{Verification record}

The census has two acceptance paths with no shared code beyond the
\texttt{geng} generator. The primary
enumeration (three programs, eight workers) emitted three canonical
JSON artifacts with embedded canonical SHA-256 hashes, listed first
below; the remaining two JSON artifacts belong to the verification
path, and the two Lean sources are hashed alongside:
\begin{center}
\small\ttfamily
\begin{tabular}{l}
t4\_support\_graph\_census.json\\
\quad 40f16a84559ace4827e366f152026f2b7868bdaed31ff9afb36184a29b48046d\\
t4\_atom\_circuit\_census.json\\
\quad 302e04ef5ff14c78cbe9dc5800ac0226e730ed0baca123585dc6469a82d66652\\
t4\_profile\_transition\_census.json\\
\quad b464682b4142a9db2396dc39ac9a0ffd8ff638aba1b9270734667c8f0a543114\\
t4\_cross\_outer\_exclusion.json\\
\quad 79db75b95e8401064f1b6159bb980ee0149f0fb3a602a607306a7f0e501a5d49\\
t4\_support\_circuit\_hit.json (Remark \ref{rem:t4:abstract})\\
\quad 5b386cd90b795bf1e6f8f174e21aa559e37c9f682e5dff373dae6bf74f3b9641\\
LiveMiddleSwapCrossOuter.lean\\
\quad 3DFF7897F65112F4F8177B84AB28F631C85BC6F50F7F5D920ED06F033B7F9275\\
R44K2TailOverlap.lean\\
\quad 12DFB92724F36BEAC1F11EBA529023F0445893CA6E60376B09B2403C0E64066A
\end{tabular}
\end{center}
Three independently written verifiers replayed the pipeline:
\texttt{verify\_t4\_support\_census.py} re-ran \texttt{geng} without
residue splitting and reproduced, via
\texttt{networkx.from\_graph6\_bytes} and NetworkX bipartitions and
distances, all \(153{,}978\) graphs, the \(34\) owner embeddings, and
the complete candidate payload
(\texttt{PASS\_INDEPENDENT\_NETWORKX\_SUPPORT\_CENSUS});
\texttt{verify\_t4\_atom\_census.py} reproduced every stage count and
the same \(576\) windows using \texttt{networkx.all\_shortest\_paths}
and NetworkX bipartite matching
(\texttt{PASS\_INDEPENDENT\_NETWORKX\_ATOM\_CENSUS}); and
\texttt{verify\_t4\_profile\_exclusion.py} reconstructed every complete
row family with its own graph6 decoder and breadth-first search,
recomputed both canonical hashes, and confirmed the forced-through
histograms of Theorem~\ref{thm:t4:census}(ii) and the absence of raw
middle swaps (\texttt{PASS\_T4\_RAW\_MIDDLE\_SWAP\_EXCLUSION}); this
last verifier was replayed once more during the audit of the session
record, with matching histograms. Neither acceptance path uses
floating-point arithmetic; the two paths use different graph6
decoders, shortest-path engines, and matching algorithms. The
\texttt{geng} counts of Table~\ref{tab:t4:geng} were additionally
re-generated for this paper. The Lean modules named above cover only
the two finite implications indicated; the census itself is not
kernel-checked, and its formal replay remains open.

\begin{remark}[Scope]\label{rem:t4:scope}
The sixteen-atom closure eliminates a single structured window. It
neither improves the general bounds on \(\bip(G)\) for triangle-free
\(G\) nor resolves the Erdos \(n^2/25\) conjecture. The neighbouring
windows \((t,k)=(5,2)\) and \((5,3)\), where
Table~\ref{tab:t4:crossover} leaves margins seven and four, are open;
the unrooted analogue of the census is infeasible there (the graph
count at \(24\) edges already exceeds \(1.9\times10^8\) at sixteen
vertices), so a rooted generation or a structural argument would be
required.
\end{remark}
